% SOURCE_STATEMENT: "Problem 20. For circular restricted three-body problem, there is a Jacobi integral. For a complete integrability, need two more integrals. Extra integral exists when m1 + m2 = 0. Is there one more integral? (Christian Marchal)"
% Preparation and final self-audit: CONTRIBUTING.md prompts C and D; source/audit artifacts retained in the local physics_statements_2026_10_11 batch.
% !TeX program = lualatex
\documentclass[11pt]{article}
\usepackage[margin=25mm]{geometry}
\usepackage{fontspec}
\setmainfont{Latin Modern Roman}
\usepackage{amsmath,amssymb,mathtools}
\usepackage{unicode-math}
\setmathfont{Latin Modern Math}
\usepackage{xurl,hyperref}
\hypersetup{colorlinks=true,urlcolor=blue}
\setcounter{secnumdepth}{0}
\setlength{\parindent}{0pt}
\setlength{\parskip}{5pt}
\setlength{\emergencystretch}{3em}
\newcommand{\sref}[2]{\hyperlink{source:#1:#2}{[S#2]}}
\newcommand{\cataloguescope}[1]{\paragraph{Recorded target.}#1}
\newcommand{\R}{\mathbb R}
\newcommand{\C}{\mathbb C}
\newcommand{\N}{\mathbb N}
\newcommand{\Z}{\mathbb Z}
\newcommand{\Q}{\mathbb Q}
\newcommand{\F}{\mathbb F}
\newcommand{\eps}{\varepsilon}
\begin{document}
% UnsolvedMath ID 20002549; source catalogue code AIM-PHYSICS-0057
\setcounter{section}{20002548}
\section{Zero-total-mass reduction and obstruction to the Euler integral}\label{20002549}
{\small UnsolvedMath ID 20002549 · Mathematical Physics}
\subsection{Definitions and mathematical statement}

\paragraph{Shared notation.}$\mathbb N=\{1,2,\ldots\}$; the standard number fields have their usual meanings. All additional conventions are stated locally.


\paragraph{Definitions and notation.}In the circular restricted three-body problem, two primaries follow prescribed circular motion and a test particle has negligible effect on them. In rotating coordinates the Jacobi constant combines kinetic energy and the effective gravitational/centrifugal potential. The Hill region is the set where the zero-velocity inequality permits motion. Its necks occur near collinear equilibria. A first integral is constant along trajectories; functional independence and Poisson commutation are needed for the usual Hamiltonian notion of complete integrability. The signed-mass condition $m_1+m_2=0$ is exceptional and cannot hold for two positive primaries. It must be interpreted in the generalized model intended by the source, not ordinary positive-mass circular dynamics.

\paragraph{Catalogue wording retained for traceability.}
In the circular restricted three-body problem, a Jacobi integral is available and the source says complete integrability requires two more integrals. It records an extra integral in the exceptional case $m_1+m_2=0$. Is there one more integral?
\paragraph{Scope and source qualification.}The signed-mass model, degree-of-freedom convention and admissible class of integrals need specification. Treating m1+m2=0 as compatible with the usual positive circular primaries would be inconsistent; the source question is preserved as an exceptional generalized setting. The catalogue title/status was paired with a machine-generated research attempt; that attempt is not peer-reviewed evidence of settlement. The statement covers the original source questions rather than only the title’s subcase.
\subsection{Short English statement}
Find whether the exceptional zero-total-mass restricted three-body model has the further first integral asked for in the source.
\subsection{Sources}
\begin{itemize}
\item[\hypertarget{source:20002549:1}{S1}] ULAM.AI and the UnsolvedMath Contributors, UnsolvedMath, record 20002549 (AIM-PHYSICS-0057). Catalogue curation under CC BY 4.0; definitions and English exposition added locally.
\url{https://huggingface.co/datasets/ulamai/UnsolvedMath}.
\item[\hypertarget{source:20002549:2}{S2}] American Institute of Mathematics, Variational Methods in Celestial Mechanics; notes by Kuo-Chang Chen, version 22 June 2003, source item 20.
\url{https://aimath.org/WWN/varcelest/varcelest.pdf}.
\end{itemize}
\label{end:20002549}
\end{document}
